\documentclass[12pt, a4paper]{article}
\usepackage[utf8]{inputenc} 
\usepackage{graphicx}     
\usepackage{float}          

\begin{document}

\subsection*{4.3. Associação entre Condições Meteorológicas e Utilização do Sistema}

Para a segunda característica selecionada na Questão 3, as condições meteorológicas (\texttt{weathersit}), a abordagem estatística requerida é diferente. Por se tratar de uma variável categórica que divide os dias em grupos distintos (1 = Bom, 2 = Nublado/Misto, 3 = Chuva/Neve Leve) a ser comparada com uma variável contínua (\texttt{total\_user}), utilizou-se a \textbf{Análise de Variância (ANOVA)}. Para a visualização da distribuição dos dados em cada categoria, optou-se pelo \textit{boxplot}.

A execução do teste ANOVA devolveu um \textit{p-value} de aproximadamente $2,13 \times 10^{-9}$. A representação visual desta distribuição encontra-se na Figura \ref{fig:boxplot_clima}.

\begin{figure}[H]
    \centering
    \caption{Associação entre condições meteorológicas e total de utilizadores diários.}
    \vspace{0.1cm}
    \includegraphics[width=0.8\textwidth]{4.3.boxplot_clima.png}
    \\ \vspace{0.1cm}
    \footnotesize{Fonte: Elaborado pelos autores (2026) com base nos dados de \texttt{2.1.total\_user.xlsx}.}
    \label{fig:boxplot_clima}
\end{figure}

\textbf{Interpretação e Comparação com Resultados Anteriores}

O \textit{p-value} obtido ($p < 0,05$) indica que existe uma diferença estatisticamente significativa na média de utilizadores entre os diferentes tipos de clima. Observando o \textit{boxplot}, é evidente que a mediana e a concentração de alugueres são consideravelmente superiores nos dias de clima limpo, sofrendo uma queda drástica nos dias de precipitação.

Esta validação matemática corrobora de forma absoluta as observações qualitativas documentadas na Questão 3. A redução do uso do sistema não é um evento aleatório, mas sim uma consequência direta e estatisticamente provada do agravamento das condições meteorológicas, consolidando o clima como um fator determinante na procura por bicicletas.

\end{document}